\begin{table}[htbp]
\centering
\begin{threeparttable}
\caption{BVAR Lag and Prior Sensitivity}
\label{tab:bvar_lag_sensitivity}
\begin{tabular}{cccccc}
\toprule
Lags ($p$) & Prior tightness ($\lambda$) & Acceptance rate (\%) & IRF sign pattern & Peak effect (bps) & 68\% band excludes 0 \\
\midrule
1 & 0.1 & 98.4 & Preserved & $-18.2$ & Yes \\
1 & 0.2 & 98.7 & Preserved & $-19.6$ & Yes \\
1 & 0.3 & 98.1 & Preserved & $-20.3$ & Yes \\
\addlinespace
2 & 0.1 & 97.9 & Preserved & $-21.4$ & Yes \\
2 & 0.2 & 98.3 & Preserved & $-22.8$ & Yes \\
2 & 0.3 & 97.6 & Preserved & $-23.1$ & Yes \\
\addlinespace
3 & 0.1 & 97.5 & Preserved & $-22.1$ & Yes \\
3 & 0.2 & 97.8 & Preserved & $-23.5$ & Yes \\
3 & 0.3 & 97.2 & Preserved & $-24.0$ & Yes \\
\addlinespace
4 & 0.1 & 97.1 & Preserved & $-21.8$ & Yes \\
4 & 0.2 & 97.4 & Preserved & $-23.2$ & Yes \\
4 & 0.3 & 97.0 & Preserved & $-23.7$ & Yes \\
\bottomrule
\end{tabular}
\begin{tablenotes}[flushleft]
\small
\item \textit{Notes:} Each row reports results from a separate BVAR estimation with the indicated lag length and Minnesota prior tightness parameter. The baseline specification uses $p=2$ and $\lambda=0.2$. ``Preserved'' indicates that the qualitative impulse response pattern (negative CPI response to an expectations shock at horizons 1--8) matches the baseline. Peak effect is the maximum absolute CPI response in basis points. The 68\% credible band column indicates whether zero is excluded at the peak response horizon. All specifications use 20,000 MCMC draws with 5,000 burn-in.
\end{tablenotes}
\end{threeparttable}
\end{table}